\documentclass[11pt,fontset=none]{ctexart}

\usepackage[a4paper,margin=1in]{geometry}
\usepackage{amsmath,amssymb,amsthm,mathtools}
\usepackage{array}
\usepackage{booktabs}
\usepackage{enumitem}
\usepackage{hyperref}
\usepackage[nameinlink,noabbrev]{cleveref}
\usepackage{xcolor}

\setCJKmainfont{STSong}
\setCJKsansfont{Heiti SC}
\setCJKmonofont{STSong}

\hypersetup{
  colorlinks=true,
  linkcolor=blue!50!black,
  citecolor=blue!50!black,
  urlcolor=blue!50!black
}

\title{exp0414\_simplexNet 工程说明文稿\\\large 现状梳理与工程化建议}
\author{Shuheng Zhang}
\date{\today}

\begin{document}

\maketitle

\begin{abstract}
本文稿面向 \texttt{memory\_lm} 仓库中的 \texttt{exp0414\_simplexNet} 子项目，给出一份中文为主的工程说明。本文不把该目录描述为完整的强化学习系统，而是将其定位为一个围绕 Simplex Memory Network (SMN) 的结构实现与实验验证原型。文稿分为三部分：其一，准确说明当前代码已经实现的能力、训练路径和测试覆盖；其二，指出当前工程封装上的缺口，尤其是 checkpoint 生命周期、本地日志、配置契约与交互确认机制；其三，给出后续工程化重构建议，包括核心模块拆分、配置组织以及基于 PySide 的训练/测试界面规划。全文只将当前仓库中可以直接核实的内容写作“现状”，并将未来方案集中放入“建议与规划”部分，避免把未实现功能误写成既成事实。
\end{abstract}

\tableofcontents

\section{项目定位}

\texttt{exp0414\_simplexNet} 的核心目标，是实现并验证一种新的前馈网络结构：Simplex Memory Network (SMN)。从当前代码路径看，该目录的主实验入口是 \texttt{run.py}，其默认工作方式是：
\[
\texttt{params.yaml} \rightarrow \texttt{build\_dataset} \rightarrow \texttt{SMNFitter.fit} \rightarrow \texttt{MLPFitter.fit (optional)} \rightarrow \texttt{plot}
\]

这一定义了它当前的真实身份：\textbf{结构原型 + 监督学习实验封装}。它不是一个现成的完整强化学习项目，也不是一个带有持久化训练状态管理的工程系统。

仓库中确实存在 \texttt{examples/smn\_for\_rl.py}，其中展示了如何把 \texttt{SMNModule} 作为 DQN 的 Q-network 或 REINFORCE 的 policy network 使用。但该文件仅是\textbf{示例接口}，并不是 \texttt{exp0414\_simplexNet} 的主实验入口，也不代表当前主流程已经转向 RL 训练。

因此，对该目录最准确的表述应为：
\begin{itemize}[leftmargin=2em]
  \item 它首先是一个新的网络结构实现项目；
  \item 它目前通过合成监督学习任务验证表达能力与训练表现；
  \item 它保留了未来接入 RL 的可能性，但 RL 仍然处在示例层，而非主线工程层。
\end{itemize}

\section{当前代码结构}

\subsection{主要脚本与职责}

当前目录的主要文件可以按职责划分为以下几类。

\begin{center}
\begin{tabular}{>{\ttfamily}p{0.28\linewidth} p{0.62\linewidth}}
\toprule
路径/模块 & 当前职责 \\
\midrule
src/config.py & 配置数据类 \texttt{Config} 与 YAML 加载逻辑。 \\
src/data.py & 合成任务定义与数据集构造，包括 1D/2D 回归任务。 \\
src/graph.py & 单纯形格点、势函数诱导有向图、拓扑层级等几何骨架。 \\
src/smn\_fitter.py & 同时包含 \texttt{SMNModule} 与 \texttt{SMNFitter}；即模型定义和高层训练封装混在同一文件。 \\
src/mlp\_fitter.py & MLP baseline 的训练封装，与 \texttt{SMNFitter} 接口对齐。 \\
src/plot.py & 可视化输出，包括散点图、loss 曲线与四图对比。 \\
run.py & 读取配置、构造数据、训练 SMN/MLP、保存图像。 \\
tests/test\_smn.py & 针对 \texttt{SMNModule} 的单元测试。 \\
examples/smn\_for\_rl.py & 将 \texttt{SMNModule} 作为 RL 函数逼近器使用的示例。 \\
\bottomrule
\end{tabular}
\end{center}

\subsection{SMNModule、SMNFitter 与 MLPFitter 的关系}

当前实现中存在两层不同抽象。

\paragraph{第一层：网络结构本身}
\texttt{SMNModule} 是一个标准的 PyTorch \texttt{nn.Module}。它负责：
\begin{itemize}[leftmargin=2em]
  \item 构造 \texttt{SimplexMemoryGraph}；
  \item 注册输入归一化所需的 \texttt{\_x\_min}/\texttt{\_x\_max} buffers；
  \item 维护边权、核心节点偏置与输出偏置；
  \item 完成按 DAG 拓扑层级推进的前向传播。
\end{itemize}

\paragraph{第二层：实验封装}
\texttt{SMNFitter} 是面向实验的高层包装器。它在 \texttt{SMNModule} 之上补了如下能力：
\begin{itemize}[leftmargin=2em]
  \item 训练循环调用；
  \item 训练/验证损失缓存；
  \item 默认 demo 数据；
  \item 预测；
  \item 图像输出。
\end{itemize}

\texttt{MLPFitter} 则提供与 \texttt{SMNFitter} 尽量相同的外层接口，以便在 \texttt{run.py} 中充当 baseline。

这一设计让实验推进很快，但也带来一个明显后果：\textbf{模型定义、训练过程、结果缓存与展示逻辑耦合在一起}。这在原型阶段可以接受，但对后续 checkpoint 恢复、GUI 调用和 RL 接入都不够理想。

\section{当前训练与测试流程}

\subsection{训练主链路}

\texttt{run.py} 的实际流程可以概括为以下步骤：
\begin{enumerate}[leftmargin=2em]
  \item 读取 \texttt{params.yaml}，构造 \texttt{Config}；
  \item 根据 \texttt{task\_name}、\texttt{n\_in}、\texttt{x\_bounds} 等设置调用 \texttt{build\_dataset}；
  \item 构造 \texttt{SMNFitter} 并执行 \texttt{fit()}；
  \item 若 \texttt{compare\_mlp=true}，则额外训练 \texttt{MLPFitter}；
  \item 调用 \texttt{plot()} 输出单模型或 SMN/MLP 对比图；
  \item 将图像保存在 \texttt{runs/<date>/<exp\_name\_timestamp>/}。
\end{enumerate}

\subsection{当前任务类型}

\texttt{src/data.py} 当前提供的是合成监督学习任务，而不是环境交互型任务。已注册任务包括：
\begin{itemize}[leftmargin=2em]
  \item 1 输入 1 输出：\texttt{sin}, \texttt{sin\_mix}, \texttt{poly\_wave}, \texttt{piecewise}
  \item 1 输入 2 输出：\texttt{sin\_cos}
  \item 2 输入 1 输出：\texttt{sin\_sum}, \texttt{sin\_product}, \texttt{quadratic}
  \item 2 输入 2 输出：\texttt{trig\_2d}
\end{itemize}

这进一步说明，当前 \texttt{0414} 的主实验仍然是函数逼近问题，而不是 MDP 训练问题。

\subsection{测试覆盖}

\texttt{tests/test\_smn.py} 当前主要覆盖 \texttt{SMNModule} 的基础行为，包括：
\begin{itemize}[leftmargin=2em]
  \item SISO/MIMO 形状正确性；
  \item 1D 输入压缩形状；
  \item 输出范围受最终 \texttt{tanh} 限制；
  \item \texttt{x\_bounds} 归一化行为；
  \item \texttt{x\_bounds} 长度不匹配时抛错；
  \item 参数数量为正；
  \item 架构字符串格式；
  \item 反向传播梯度可流动；
  \item 若干 \((n,m)\) 组合的 smoke test；
  \item \texttt{state\_dict} save/load roundtrip。
\end{itemize}

这意味着当前测试重点是\textbf{结构级正确性}，而不是完整训练系统的状态管理、结果重现或多轮实验恢复。

\section{当前已具备的能力}

基于当前代码，可以确认该项目已经具备以下能力。

\subsection{SMN 结构实现已完整落地}

\texttt{graph.py} 已经实现：
\begin{itemize}[leftmargin=2em]
  \item 单纯形格点枚举；
  \item 输入面 \texttt{F\_in}、输出面 \texttt{F\_out}、共享中间面 \texttt{F\_mid}；
  \item 势函数驱动的有向最近邻边；
  \item 前驱/后继表；
  \item 拓扑层级计算与验证。
\end{itemize}

\texttt{SMNModule} 则将这一几何骨架变为实际可训练网络。

\subsection{支持多种输入输出形态}

当前实现不仅支持 SISO，也支持更一般的 MIMO 场景。输入归一化通过每通道 \texttt{x\_bounds} 完成；输出维度由 \texttt{n\_out} 控制。这一点已经在配置、数据生成和测试中得到体现。

\subsection{提供 MLP baseline 对照}

\texttt{MLPFitter} 被设计为与 \texttt{SMNFitter} 拥有近似一致的接口，因此可以在同一训练脚本中直接进行结构比较。当前 \texttt{run.py} 默认就是以 SMN 为主角、MLP 为对照。

\subsection{支持 state\_dict 级别的模型序列化接口}

作为标准 \texttt{nn.Module}，\texttt{SMNModule} 已支持 \texttt{state\_dict()} 与 \texttt{load\_state\_dict()}。测试中也已有 roundtrip 检验。这意味着\textbf{模型参数本身是可保存、可恢复的}。

不过，需要强调的是：当前支持的是\textbf{接口层面的序列化能力}，而不是完整的\textbf{本地 checkpoint 生命周期管理系统}。这一点将在后文说明。

\subsection{存在 RL 使用示例，但仅处于示例层}

\texttt{examples/smn\_for\_rl.py} 中包含 DQN 与 REINFORCE 风格的样例，说明 \texttt{SMNModule} 可以作为状态到 Q 值、或状态到策略 logits 的映射器使用。这一层证明了该结构与 RL 的接口是可行的。

但当前目录并未提供一个以 RL 为主线的正式训练入口，因此不应把该示例误判为该目录已经实现了完整 RL 框架。

\section{当前缺口与封装问题}

\subsection{缺少本地 checkpoint 生命周期管理}

当前训练逻辑在 \texttt{smn\_fitter.py} 与 \texttt{mlp\_fitter.py} 中都维护了一个内存中的 \texttt{best\_state}，训练完成后会把最优参数重新加载回模块。但这一过程只发生在内存里，\textbf{不会自动写入磁盘 checkpoint}。

因此，当前系统不存在以下工程能力：
\begin{itemize}[leftmargin=2em]
  \item 自动发现已有本地权重；
  \item 兼容时优先加载本地权重；
  \item 不兼容时展示差异并由用户确认；
  \item 初始化新模型后自动建立可追踪的 checkpoint 历史；
  \item 中断后继续训练。
\end{itemize}

\subsection{缺少配置兼容性比较与替换确认}

从当前代码看，训练系统不会检查“本地已有参数是否与当前配置相容”，也不会在发生初始化、替换或覆盖前弹出确认。对研究型实验来说，这意味着已有结果很容易被忽略；对未来 RL 训练来说，这则是更高风险的问题。

后续如果要把该目录升级为可恢复训练系统，至少应当支持：
\begin{itemize}[leftmargin=2em]
  \item checkpoint 元信息记录；
  \item 配置差异比较；
  \item \texttt{state\_dict} key/shape 差异报告；
  \item 初始化或替换前的确认弹窗；
  \item 决策结果写入日志。
\end{itemize}

\subsection{缺少持续追加的本地运行日志}

当前目录会在 \texttt{runs/} 下输出图片，但缺少一份持续追加的本地日志文件，用于记录每次运行的配置、加载来源、差异检查结果、训练过程与保存动作。这使得实验历史难以追溯，也不利于后续排查“哪次训练用了哪个模型、在哪一步被替换”。

\subsection{配置层与执行层已有漂移}

\texttt{Config} 与 \texttt{params.yaml} 中存在一些当前未真正接入主流程的字段，例如：
\begin{itemize}[leftmargin=2em]
  \item \texttt{model\_type}
  \item \texttt{window\_width}
  \item \texttt{window\_hold}
\end{itemize}

\texttt{run.py} 当前仍然固定训练 SMN，并在需要时附带训练 MLP baseline，而不是根据 \texttt{model\_type} 选择主模型。同样，窗口训练参数已经出现在配置中，但未实际驱动训练主链路。对工程系统而言，这种“配置表面大于实际实现”的状态容易误导使用者。

\section{工程化建议}

\subsection{建议的核心模块拆分}

若要把 \texttt{exp0414\_simplexNet} 从“实验原型”推进为“可恢复、可审计、可扩展”的工程系统，更清晰的主干应当拆成四个核心模块：
\begin{description}[leftmargin=2em,style=nextline]
  \item[\texttt{problem}] 定义问题本身。当前阶段可承载监督学习任务定义；若未来接入 RL，则承载 MDP、状态动作空间和采样接口。
  \item[\texttt{model}] 定义网络结构本身。SMN 的几何图与 \texttt{SMNModule} 都应归于这一层。
  \item[\texttt{train}] 定义训练流程，包括 optimizer、loss、checkpoint、日志和恢复逻辑。
  \item[\texttt{test/eval}] 定义单元测试、评估流程、结果分析与回归验证。
\end{description}

相应地，\texttt{plot}、\texttt{logs}、\texttt{checkpoints}、\texttt{docs} 则应作为辅助层，而不应反向侵入训练主流程。

\subsection{checkpoint 与确认机制建议}

后续 checkpoint 机制应遵循如下保守原则：
\begin{enumerate}[leftmargin=2em]
  \item 启动时先检查本地已有权重；
  \item 若本地权重与当前配置兼容，则优先加载，而不是重新初始化；
  \item 若不兼容，则明确展示差异，说明不兼容发生在哪些字段、哪些参数键或哪些张量形状上；
  \item 只有在用户确认后，才允许初始化新模型或替换旧模型；
  \item 所有初始化、替换、加载与保存动作都写入本地追加日志。
\end{enumerate}

这套机制的目标不是“自动化越多越好”，而是“\textbf{对已有实验成果尽量保守，对覆盖行为尽量显式确认}”。

\subsection{配置组织建议}

后续配置不宜继续只依赖一份混合式 \texttt{params.yaml}。更稳妥的方案是拆成三类配置：
\begin{itemize}[leftmargin=2em]
  \item \texttt{model.yaml}：只描述模型结构与输入输出边界；
  \item \texttt{train.yaml}：只描述训练超参数、checkpoint 策略、日志策略；
  \item \texttt{eval.yaml}：只描述测试/评估所需的模型来源、任务、指标与输出选项。
\end{itemize}

其中，模型结构字段应在训练与测试之间共享定义，以避免配置契约漂移。

\section{配置与 GUI 规划}

\subsection{训练与测试界面应独立，但风格统一}

如果后续要给该项目补上桌面交互层，推荐采用 PySide 构建两个独立窗口：
\begin{itemize}[leftmargin=2em]
  \item \textbf{训练窗口}：用于选择/比对 checkpoint、编辑训练配置、启动训练、查看日志和过程曲线；
  \item \textbf{测试/评估窗口}：用于加载模型、选择任务或评估数据、运行评估并查看结果。
\end{itemize}

这两个界面应当共享统一的视觉系统，例如字体、色彩、卡片布局、确认弹窗样式和日志面板风格，但不应合并为一个职责过于混杂的“大一统窗口”。

\subsection{GUI 的职责边界}

需要强调，GUI 只能是交互壳层，而不应承载核心训练逻辑。较合理的职责分工是：
\begin{itemize}[leftmargin=2em]
  \item GUI 负责读取/写入 YAML、展示差异、弹出确认框、启动任务、展示日志与结果；
  \item 后端模块负责真实的训练、恢复、保存、评估与错误处理。
\end{itemize}

只有这样，后续 CLI 批处理、自动实验脚本或其他前端入口才不会被 GUI 实现反向绑死。

\section{结论}

\texttt{exp0414\_simplexNet} 当前的价值，不在于它已经成为一个完整的训练平台，而在于它已经把一种新型网络结构 SMN 实实在在地编码成了可运行、可训练、可测试、可对照的原型系统。它完成了最关键的一步：证明“结构”本身是可以落地的。

但若要继续向长期使用的实验系统推进，当前还缺少若干工程层能力，尤其是 checkpoint 生命周期、差异比较、确认机制、追加日志、配置契约与 GUI 分层。这些问题并不否定当前目录的研究价值，恰恰说明它正处在一个清晰的下一阶段入口上：
\[
\text{从结构原型与实验封装} \quad \longrightarrow \quad \text{可恢复、可审计、可交互的工程系统}
\]

因此，本文稿的基本判断是：\texttt{0414} 已经是一个很好的网络结构研究原型，但它的下一步不应只是“多跑几组实验”，而应是把训练、测试、配置、持久化和交互层系统性地补齐。

\end{document}
